% !TeX program = pdflatex
\documentclass[11pt,a4paper]{amsart}

\newcommand{\notelanguage}{finnish}
% Shared preamble for course notes and exercise sets.

\usepackage[T1]{fontenc}
\usepackage{lmodern}
\providecommand{\notelanguage}{english}
\usepackage[finnish,english]{babel}
\AtBeginDocument{\expandafter\selectlanguage\expandafter{\notelanguage}}

\newcommand{\notesfinnishlanguage}{finnish}
\ifx\notelanguage\notesfinnishlanguage
\newcommand{\notestheoremname}{Lause}
\newcommand{\noteslemmaname}{Lemma}
\newcommand{\notespropositionname}{Propositio}
\newcommand{\notescorollaryname}{Korollaari}
\newcommand{\notesdefinitionname}{Määritelmä}
\newcommand{\notesexamplename}{Esimerkki}
\newcommand{\notesproblemname}{Tehtävä}
\newcommand{\notesexercisename}{Tehtävä}
\newcommand{\notesremarkname}{Huomautus}
\newcommand{\notessolutionname}{Ratkaisu}
\else
\newcommand{\notestheoremname}{Theorem}
\newcommand{\noteslemmaname}{Lemma}
\newcommand{\notespropositionname}{Proposition}
\newcommand{\notescorollaryname}{Corollary}
\newcommand{\notesdefinitionname}{Definition}
\newcommand{\notesexamplename}{Example}
\newcommand{\notesproblemname}{Problem}
\newcommand{\notesexercisename}{Exercise}
\newcommand{\notesremarkname}{Remark}
\newcommand{\notessolutionname}{Solution}
\fi

\usepackage[a4paper,margin=30mm]{geometry}
\usepackage{microtype}
\usepackage{graphicx}
\usepackage{enumitem}

\newcommand{\exerciseheading}[2]{%
  \par\noindent
  {\large\bfseries #1\par}
  \vspace{0.1em}
  \noindent{\large #2\par}
  \vspace{1.5em}
}

\usepackage{amsmath}
\usepackage{amssymb}
\usepackage{amsthm}
\usepackage{mathtools}
\usepackage{aliascnt}

% Row operations beside the affected row, e.g. \rowop{+2R_1}.
\newcommand{\rowop}[1]{\mathrel{\big|}\,#1}

% Linear systems: the argument gives the number of alignment pairs.
\newenvironment{linearsystem}[1]
{\left\{
  \begin{alignedat}{#1}}
    {
  \end{alignedat}\right.}

  \usepackage[hidelinks,pdfusetitle]{hyperref}

  \numberwithin{equation}{section}
  \setlist{itemsep=0.25em,topsep=0.5em}

  \theoremstyle{plain}
  \newtheorem{theorem}{\notestheoremname}[section]
  \newaliascnt{lemma}{theorem}
  \newtheorem{lemma}[lemma]{\noteslemmaname}
  \aliascntresetthe{lemma}
  \newaliascnt{proposition}{theorem}
  \newtheorem{proposition}[proposition]{\notespropositionname}
  \aliascntresetthe{proposition}
  \newaliascnt{corollary}{theorem}
  \newtheorem{corollary}[corollary]{\notescorollaryname}
  \aliascntresetthe{corollary}

  \theoremstyle{definition}
  \newaliascnt{definition}{theorem}
  \newtheorem{definition}[definition]{\notesdefinitionname}
  \aliascntresetthe{definition}
  \newaliascnt{example}{theorem}
  \newtheorem{example}[example]{\notesexamplename}
  \aliascntresetthe{example}
  \newaliascnt{problem}{theorem}
  \newtheorem{problem}[problem]{\notesproblemname}
  \aliascntresetthe{problem}
  \newtheorem{exercise}{\notesexercisename}

  \theoremstyle{remark}
  \newaliascnt{remark}{theorem}
  \newtheorem{remark}[remark]{\notesremarkname}
  \aliascntresetthe{remark}

  \newenvironment{solution}
  {
  \begin{proof}[\notessolutionname]}
    {
  \end{proof}}

  \usepackage[nameinlink,noabbrev]{cleveref}

  \crefname{theorem}{theorem}{theorems}
  \crefname{lemma}{lemma}{lemmas}
  \crefname{proposition}{proposition}{propositions}
  \crefname{corollary}{corollary}{corollaries}
  \crefname{definition}{definition}{definitions}
  \crefname{example}{example}{examples}
  \crefname{problem}{problem}{problems}
  \crefname{exercise}{exercise}{exercises}
  \crefname{remark}{remark}{remarks}

  \ifx\notelanguage\notesfinnishlanguage
  \crefname{theorem}{lause}{lauseet}
  \crefname{lemma}{lemma}{lemmat}
  \crefname{proposition}{propositio}{propositiot}
  \crefname{corollary}{korollaari}{korollaarit}
  \crefname{definition}{määritelmä}{määritelmät}
  \crefname{example}{esimerkki}{esimerkit}
  \crefname{problem}{tehtävä}{tehtävät}
  \crefname{exercise}{tehtävä}{tehtävät}
  \crefname{remark}{huomautus}{huomautukset}
  \fi

  \newcommand{\NN}{\mathbb{N}}
  \newcommand{\NNone}{\mathbb{N}_1}
  \newcommand{\ZZ}{\mathbb{Z}}
  \newcommand{\QQ}{\mathbb{Q}}
  \newcommand{\RR}{\mathbb{R}}
  \newcommand{\CC}{\mathbb{C}}
  \newcommand{\PP}{\mathbb{P}}

  \DeclareMathOperator{\im}{im}
  \DeclareMathOperator{\rank}{rank}
  \DeclareMathOperator{\syt}{syt}
  \DeclareMathOperator{\pyj}{pyj}
  \DeclareMathOperator{\dom}{dom}
  \DeclarePairedDelimiter{\abs}{\lvert}{\rvert}
  \DeclarePairedDelimiter{\norm}{\lVert}{\rVert}
  \newcommand{\set}[1]{\left\lbrace#1\right\rbrace}
  \newcommand{\maxset}[1]{\max\set{#1}}
  \newcommand{\minset}[1]{\min\set{#1}}
  \newcommand{\powerset}[1]{\mathcal{P}\!\left(#1\right)}


\title{Lineaarialgebra ja matriisilaskenta I\\Harjoituksen 2 ratkaisut}
\author{}
\date{}

\begin{document}

\exerciseheading{Lineaarialgebra ja matriisilaskenta I}{Harjoituksen
2 ratkaisut}

\begin{exercise}
  Olkoot
  \[
    A =
    \begin{bmatrix}
      2 & -1 & -4 \\
      -5 & 3 & 0 \\
      -2 & -3 & 1
    \end{bmatrix}
    \qquad\text{ja}\qquad
    B =
    \begin{bmatrix}
      1 & 7 & -1 \\
      -2 & -2 & 2 \\
      5 & 2 & -3
    \end{bmatrix}.
  \]
  Laske \(A+B\), \(B-2A\) ja \(AB\). Onko voimassa \(AB=BA\)?
\end{exercise}

\begin{solution}
  Laskemalla saadaan
  \[
    \begin{aligned}
      A + B &=
      \begin{bmatrix}
        2 + 1 & -1 + 7 & -4 - 1 \\
        -5 - 2 & 3 - 2 & 0 + 2 \\
        -2 + 5 & -3 + 2 & 1 - 3
      \end{bmatrix} \\
      &=
      \begin{bmatrix}
        3 & 6 & -5 \\
        -7 & 1 & 2 \\
        3 & -1 & -2
      \end{bmatrix}.
    \end{aligned}
  \]
  Samoin
  \[
    \begin{aligned}
      B - 2A &=
      \begin{bmatrix}
        1 - 4 & 7 + 2 & -1 + 8 \\
        -2 + 10 & -2 - 6 & 2 - 0 \\
        5 + 4 & 2 + 6 & -3 - 2
      \end{bmatrix} \\
      &=
      \begin{bmatrix}
        -3 & 9 & 7 \\
        8 & -8 & 2 \\
        9 & 8 & -5
      \end{bmatrix}.
    \end{aligned}
  \]
  Lasketaan seuraavaksi matriisitulo \(AB\) sarakettain. Merkintä
  \(B_{\bullet j}\) tarkoittaa matriisin \(B\) saraketta \(j\), joten
  \((AB)_{\bullet j} = A B_{\bullet j}\). Saadaan
  \[
    \begin{aligned}
      AB_{\bullet 1} &=
      \begin{bmatrix}
        2\cdot 1 + (-1)\cdot(-2) + (-4)\cdot 5 \\
        (-5)\cdot 1 + 3\cdot(-2) + 0\cdot 5 \\
        (-2)\cdot 1 + (-3)\cdot(-2) + 1\cdot 5
      \end{bmatrix}
      &=
      \begin{bmatrix}
        -16 \\
        -11 \\
        9
      \end{bmatrix}, \\
      AB_{\bullet 2} &=
      \begin{bmatrix}
        2\cdot 7 + (-1)\cdot(-2) + (-4)\cdot 2 \\
        (-5)\cdot 7 + 3\cdot(-2) + 0\cdot 2 \\
        (-2)\cdot 7 + (-3)\cdot(-2) + 1\cdot 2
      \end{bmatrix}
      &=
      \begin{bmatrix}
        8 \\
        -41 \\
        -6
      \end{bmatrix}, \\
      AB_{\bullet 3} &=
      \begin{bmatrix}
        2\cdot(-1) + (-1)\cdot 2 + (-4)\cdot(-3) \\
        (-5)\cdot(-1) + 3\cdot 2 + 0\cdot(-3) \\
        (-2)\cdot(-1) + (-3)\cdot 2 + 1\cdot(-3)
      \end{bmatrix}
      &=
      \begin{bmatrix}
        8 \\
        11 \\
        -7
      \end{bmatrix}.
    \end{aligned}
  \]
  Täten
  \[
    AB =
    \begin{bmatrix}
      (AB)_{\bullet 1} & (AB)_{\bullet 2} & (AB)_{\bullet 3}
    \end{bmatrix} =
    \begin{bmatrix}
      -16 & 8 & 8 \\
      -11 & -41 & 11 \\
      9 & -6 & -7
    \end{bmatrix}.
  \]
  Huomataan, että
  \[
    (BA)_{13} = 1\cdot(-4) + 7\cdot 0 + (-1)\cdot 1
    = -5 \neq 8 = (AB)_{13}.
  \]
  Täten \(AB \neq BA\).
\end{solution}

\begin{exercise}
  Olkoon \(A =
    \begin{bmatrix} 1 & 1 \\ 1 & 1
  \end{bmatrix}\).
  Osoita, että
  \[
    A^n =
    \begin{bmatrix}
      2^{n-1} & 2^{n-1} \\
      2^{n-1} & 2^{n-1}
    \end{bmatrix},
    \qquad\text{kun } n \in \ZZ_+.
  \]
  \textit{Vihje:} Induktio.
\end{exercise}

\begin{solution}
  Koska \(2^0 = 1\), saadaan
  \[
    A^1 = A =
    \begin{bmatrix}
      1 & 1 \\
      1 & 1
    \end{bmatrix}
    =
    \begin{bmatrix}
      2^{1 - 1} & 2^{1 - 1} \\
      2^{1 - 1} & 2^{1 - 1}
    \end{bmatrix}.
  \]
  Väite siis pätee, kun \(n = 1\). Olkoon \(n \in \ZZ_{+}\) ja
  oletetaan, että
  \[
    A^n =
    \begin{bmatrix}
      2^{n - 1} & 2^{n - 1} \\
      2^{n - 1} & 2^{n - 1}
    \end{bmatrix}.
  \]
  Nyt käyttämällä induktio-oletusta ja laskemalla saadaan
  \[
    \begin{aligned}
      A^{n + 1} &= AA^n \\
      &=
      \begin{bmatrix}
        1 \cdot 2^{n - 1} + 1 \cdot 2^{n - 1} & 1 \cdot 2^{n - 1} + 1
        \cdot 2^{n - 1} \\
        1 \cdot 2^{n - 1} + 1 \cdot 2^{n - 1} & 1 \cdot 2^{n - 1} + 1
        \cdot 2^{n - 1}
      \end{bmatrix} \\
      &=
      \begin{bmatrix}
        2 \cdot 2^{n - 1} & 2 \cdot 2^{n - 1} \\
        2 \cdot 2^{n - 1} & 2 \cdot 2^{n - 1}
      \end{bmatrix} \\
      &=
      \begin{bmatrix}
        2^n & 2^n \\
        2^n & 2^n
      \end{bmatrix} \\
      &=
      \begin{bmatrix}
        2^{(n + 1) - 1} & 2^{(n + 1) - 1} \\
        2^{(n + 1) - 1} & 2^{(n + 1) - 1}
      \end{bmatrix}.
    \end{aligned}
  \]
  Täten väite pätee luvulla \(n + 1\). Induktioperiaatteen nojalla väite
  pätee kaikilla \(n \in \ZZ_{+}\).
\end{solution}

\begin{exercise}
  Olkoot
  \[
    A =
    \begin{bmatrix}
      1 & 1 & -2 \\
      2 & 3 & 4
    \end{bmatrix}
    \qquad\text{ja}\qquad
    b =
    \begin{bmatrix} 5 \\ 2
    \end{bmatrix}.
  \]
  Etsi matriisiyhtälön \(Ax=b\) ratkaisut \(x \in \RR^{3\times 1}\)
  yhtälöryhmän \([A \mid b]\) avulla.
\end{exercise}

\begin{solution}
  Muokataan yhtälöryhmän liitetty matriisi supistettuun porrasmuotoon
  rivioperaatioilla:
  \[
    \begin{aligned}
      [A \mid b] &={}
      \left[
        \begin{array}{ccc|r}
          1 & 1 & -2 & 5 \\
          2 & 3 & 4 & 2
        \end{array}
      \right]
      \quad
      \begin{array}{@{}l@{}}
        {} \\
        \smash{\rowop{-2R_1}}
      \end{array}
      \\[1ex]
      &\sim{}
      \left[
        \begin{array}{ccc|r}
          1 & 1 & -2 & 5 \\
          0 & 1 & 8 & -8
        \end{array}
      \right]
      \quad
      \begin{array}{@{}l@{}}
        \smash{\rowop{-R_2}} \\
        {}
      \end{array}
      \\[1ex]
      &\sim{}
      \left[
        \begin{array}{ccc|r}
          1 & 0 & -10 & 13 \\
          0 & 1 & 8 & -8
        \end{array}
      \right].
    \end{aligned}
  \]
  Täten \(x_1 = 10x_3 + 13\) ja \(x_2 = -8x_3 - 8\).
  Muuttuja \(x_3\) on vapaa. Valitsemalla \(x_3 = t \in \RR\) saadaan
  \[
    x =
    \begin{bmatrix} x_1 \\ x_2 \\ x_3
    \end{bmatrix}
    =
    \begin{bmatrix} 10t+13 \\ -8t-8 \\ t
    \end{bmatrix}
    =
    \begin{bmatrix} 13 \\ -8 \\ 0
    \end{bmatrix}
    + t
    \begin{bmatrix} 10 \\ -8 \\ 1
    \end{bmatrix}.
  \]
  Koska rivioperaatiot säilyttävät ratkaisujoukon, matriisiyhtälön
  \(Ax = b\) ratkaisujoukko on
  \[
    \left\{
      \begin{bmatrix} 13 \\ -8 \\ 0
      \end{bmatrix}
      + t
      \begin{bmatrix} 10 \\ -8 \\ 1
      \end{bmatrix}
      \colon t \in \RR
    \right\}.
  \]
\end{solution}

\begin{exercise}
  Olkoon \(A \in \RR^{m\times n}\) matriisi. Määritellään kuvaus
  \(L_A\colon \RR^{n\times 1} \to \RR^{m\times 1}\) ehdolla
  \[
    L_A(x) = Ax, \qquad x \in \RR^{n\times 1}.
  \]
  Osoita eksplisiittisellä matriisitulojen laskulla, että kuvaukselle
  \(L_A\) on voimassa
  \[
    L_A(x+y) = L_A(x)+L_A(y)
    \qquad\text{ja}\qquad
    L_A(\lambda x) = \lambda L_A(x)
  \]
  kaikilla \(x,y \in \RR^{n\times 1}\) ja \(\lambda \in \RR\).
  (Tehtävän tarkoituksena on tarkistaa monisteen Lemman 2.3.15
  todistuksen laskut kirjoittamalla matriisitulot auki.)
\end{exercise}

\begin{solution}
  Olkoot \(x,y \in \RR^{n\times 1}\) ja \(\lambda \in \RR\)
  mielivaltaisia. Kirjoitetaan
  \[
    A =
    \begin{bmatrix}
      a_{11} & \cdots & a_{1n} \\
      \vdots & \ddots & \vdots \\
      a_{m1} & \cdots & a_{mn}
    \end{bmatrix},
    \qquad
    x =
    \begin{bmatrix} x_1 \\ \vdots \\ x_n
    \end{bmatrix},
    \qquad
    y =
    \begin{bmatrix} y_1 \\ \vdots \\ y_n
    \end{bmatrix}.
  \]
  Matriisitulon määritelmän sekä reaalilukujen osittelulain ja
  yhteenlaskun laskusääntöjen nojalla saadaan
  \[
    \begin{aligned}
      L_A(x+y)
      &= A(x+y)
      = A
      \begin{bmatrix} x_1+y_1 \\ \vdots \\ x_n+y_n
      \end{bmatrix}
      \\[1ex]
      &=
      \begin{bmatrix}
        a_{11}(x_1+y_1)+\cdots+a_{1n}(x_n+y_n) \\
        \vdots \\
        a_{m1}(x_1+y_1)+\cdots+a_{mn}(x_n+y_n)
      \end{bmatrix}
      \\[1ex]
      &=
      \begin{bmatrix}
        (a_{11}x_1+\cdots+a_{1n}x_n)
        +(a_{11}y_1+\cdots+a_{1n}y_n) \\
        \vdots \\
        (a_{m1}x_1+\cdots+a_{mn}x_n)
        +(a_{m1}y_1+\cdots+a_{mn}y_n)
      \end{bmatrix}
      \\[1ex]
      &=
      \begin{bmatrix}
        a_{11}x_1+\cdots+a_{1n}x_n \\
        \vdots \\
        a_{m1}x_1+\cdots+a_{mn}x_n
      \end{bmatrix}
      +
      \begin{bmatrix}
        a_{11}y_1+\cdots+a_{1n}y_n \\
        \vdots \\
        a_{m1}y_1+\cdots+a_{mn}y_n
      \end{bmatrix}
      \\[1ex]
      &= Ax+Ay = L_A(x)+L_A(y).
    \end{aligned}
  \]
  Vastaavasti reaalilukujen kertolaskun laskusääntöjen ja osittelulain
  nojalla
  \[
    \begin{aligned}
      L_A(\lambda x)
      &= A(\lambda x)
      = A
      \begin{bmatrix} \lambda x_1 \\ \vdots \\ \lambda x_n
      \end{bmatrix}
      \\[1ex]
      &=
      \begin{bmatrix}
        a_{11}(\lambda x_1)+\cdots+a_{1n}(\lambda x_n) \\
        \vdots \\
        a_{m1}(\lambda x_1)+\cdots+a_{mn}(\lambda x_n)
      \end{bmatrix}
      \\[1ex]
      &= \lambda
      \begin{bmatrix}
        a_{11}x_1+\cdots+a_{1n}x_n \\
        \vdots \\
        a_{m1}x_1+\cdots+a_{mn}x_n
      \end{bmatrix}
      \\[1ex]
      &= \lambda(Ax) = \lambda L_A(x).
    \end{aligned}
  \]
  Molemmat yhtälöt pätevät siis kaikilla \(x,y \in \RR^{n\times 1}\)
  ja \(\lambda \in \RR\), joten \(L_A\) on lineaarinen.
\end{solution}

\begin{exercise}
  Tutki, onko matriisi
  \[
    \text{(a)}\;
    \begin{bmatrix}
      -1 & 3 & -4 \\
      2 & 4 & 1 \\
      -4 & 2 & -9
    \end{bmatrix},
    \qquad
    \text{(b)}\;
    \begin{bmatrix}
      \frac15 & \frac15 & -\frac25 \\
      \frac15 & \frac15 & \frac1{10} \\
      \frac15 & -\frac45 & \frac1{10}
    \end{bmatrix}
  \]
  kääntyvä, ja myönteisessä tapauksessa määritä matriisin käänteismatriisi.
\end{exercise}

\newpage
\begin{solution}
  \leavevmode\par
  \begin{enumerate}[label=(\alph*), leftmargin=*]
    \item Merkitään annettua matriisia kirjaimella \(A\).
      Neliömatriisien kääntyvyyslauseen (Lause 2.5.1)
      mukaan \(A\) on kääntyvä täsmälleen silloin, kun yhtälön
      \(Ax=0\) ainoa ratkaisu on nollavektori. Etsitään nollasta
      poikkeava ratkaisu. Ensimmäiseltä riviltä saadaan
      \(-x_1+3x_2-4x_3=0\), eli \(x_1=3x_2-4x_3\).
      Sijoittamalla tämä kahteen muuhun yhtälöön saadaan
      \[
        \begin{aligned}
          2(3x_2-4x_3)+4x_2+x_3 &= 10x_2-7x_3=0, \\
          -4(3x_2-4x_3)+2x_2-9x_3 &= -10x_2+7x_3=0.
        \end{aligned}
      \]
      Molemmat ehdot toteutuvat, kun \(x_2=\frac7{10}x_3\).
      Valitaan \(x_3=10\), jolloin \(x_2=7\) ja
      \(x_1=3\cdot7-4\cdot10=-19\). Tarkistetaan vielä sijoittamalla:
      \[
        A
        \begin{bmatrix} -19 \\ 7 \\ 10
        \end{bmatrix}
        =
        \begin{bmatrix}
          19+21-40 \\
          -38+28+10 \\
          76+14-90
        \end{bmatrix}
        =
        \begin{bmatrix} 0 \\ 0 \\ 0
        \end{bmatrix}.
      \]
      Yhtälöllä \(Ax=0\) on siis nollasta poikkeava ratkaisu, joten
      kääntyvyyslauseen nojalla \(A\) ei ole kääntyvä.

    \item Merkitään tämän kohdan matriisia kirjaimella \(B\).
      Muokataan liitettyä matriisia \([B\mid I]\) rivioperaatioilla:
      \[
        \begin{aligned}
          [B\mid I] &={}
          \left[
            \begin{array}{rrr|rrr}
              \frac15 & \frac15 & -\frac25 & 1 & 0 & 0 \\
              \frac15 & \frac15 & \frac1{10} & 0 & 1 & 0 \\
              \frac15 & -\frac45 & \frac1{10} & 0 & 0 & 1
            \end{array}
          \right]
          \quad
          \begin{array}{@{}l@{}}
            \vphantom{\frac15}\smash{\rowop{\cdot 5}} \\
            \vphantom{\frac1{10}}\smash{\rowop{\cdot 10}} \\
            \vphantom{\frac1{10}}\smash{\rowop{\cdot 10}}
          \end{array}
          \\[1ex]
          &\sim{}
          \left[
            \begin{array}{rrr|rrr}
              1 & 1 & -2 & 5 & 0 & 0 \\
              2 & 2 & 1 & 0 & 10 & 0 \\
              2 & -8 & 1 & 0 & 0 & 10
            \end{array}
          \right]
          \quad
          \begin{array}{@{}l@{}}
            {} \\
            \smash{\rowop{-2R_1}} \\
            \smash{\rowop{-2R_1}}
          \end{array}
          \\[1ex]
          &\sim{}
          \left[
            \begin{array}{rrr|rrr}
              1 & 1 & -2 & 5 & 0 & 0 \\
              0 & 0 & 5 & -10 & 10 & 0 \\
              0 & -10 & 5 & -10 & 0 & 10
            \end{array}
          \right]
          \quad
          \begin{array}{@{}l@{}}
            {} \\
            \smash{\rowop{R_2\leftrightarrow R_3}} \\
            {}
          \end{array}
          \\[1ex]
          &\sim{}
          \left[
            \begin{array}{rrr|rrr}
              1 & 1 & -2 & 5 & 0 & 0 \\
              0 & -10 & 5 & -10 & 0 & 10 \\
              0 & 0 & 5 & -10 & 10 & 0
            \end{array}
          \right]
          \quad
          \begin{array}{@{}l@{}}
            {} \\
            \smash{\rowop{\cdot(-1/10)}} \\
            \smash{\rowop{\cdot(1/5)}}
          \end{array}
        \end{aligned}
      \]
      Normalisoidaan johtavat alkiot ja poistetaan niiden yläpuoliset
      alkiot:
      \[
        \begin{aligned}
          &\sim{}
          \left[
            \begin{array}{rrr|rrr}
              1 & 1 & -2 & 5 & 0 & 0 \\
              0 & 1 & -\frac12 & 1 & 0 & -1 \\
              0 & 0 & 1 & -2 & 2 & 0
            \end{array}
          \right]
          \quad
          \begin{array}{@{}l@{}}
            \smash{\rowop{-R_2}} \\
            \vphantom{\frac12} \\
            {}
          \end{array}
          \\[1ex]
          &\sim{}
          \left[
            \begin{array}{rrr|rrr}
              1 & 0 & -\frac32 & 4 & 0 & 1 \\
              0 & 1 & -\frac12 & 1 & 0 & -1 \\
              0 & 0 & 1 & -2 & 2 & 0
            \end{array}
          \right]
          \quad
          \begin{array}{@{}l@{}}
            \vphantom{\frac32}\smash{\rowop{+(3/2)R_3}} \\
            \vphantom{\frac12}\smash{\rowop{+(1/2)R_3}} \\
            {}
          \end{array}
          \\[1ex]
          &\sim{}
          \left[
            \begin{array}{rrr|rrr}
              1 & 0 & 0 & 1 & 3 & 1 \\
              0 & 1 & 0 & 0 & 1 & -1 \\
              0 & 0 & 1 & -2 & 2 & 0
            \end{array}
          \right].
        \end{aligned}
      \]
      Oikeasta lohkosta saadaan käänteismatriisiehdokas
      \[
        C =
        \begin{bmatrix}
          1 & 3 & 1 \\
          0 & 1 & -1 \\
          -2 & 2 & 0
        \end{bmatrix}.
      \]
      Suoralla matriisikertolaskulla saadaan
      \[
        BC=CB=
        \begin{bmatrix}
          1 & 0 & 0 \\
          0 & 1 & 0 \\
          0 & 0 & 1
        \end{bmatrix}.
      \]
      Siis \(BC=CB=I\), joten käänteismatriisin määritelmän nojalla
      \(B\) on kääntyvä ja \(B^{-1}=C\).
  \end{enumerate}
\end{solution}

\begin{exercise}
  Olkoon
  \[
    A =
    \begin{bmatrix}
      2 & 0 & 3 & 2 \\
      1 & 1 & -1 & 0 \\
      0 & 0 & 1 & 0 \\
      1 & 0 & 0 & 1
    \end{bmatrix} \in \RR^{4\times 4}.
  \]
  Selvitä neliömatriisien kääntyvyyslauseen (Lause 2.5.1) avulla, onko
  matriisi \(A\) kääntyvä.
\end{exercise}

\begin{solution}
  Lauseen 2.5.1 mukaan \(A\) on kääntyvä täsmälleen silloin,
  kun yhtälön \(Ax=0\) ainoa ratkaisu on nollavektori.
  Tutkitaan homogeenista yhtälöryhmää rivioperaatioilla:
  \[
    \begin{aligned}
      [A\mid 0] &={}
      \left[
        \begin{array}{rrrr|r}
          2 & 0 & 3 & 2 & 0 \\
          1 & 1 & -1 & 0 & 0 \\
          0 & 0 & 1 & 0 & 0 \\
          1 & 0 & 0 & 1 & 0
        \end{array}
      \right]
      \quad
      \begin{array}{@{}l@{}}
        \smash{\rowop{-2R_4}} \\
        \smash{\rowop{-R_4}} \\
        {} \\
        {}
      \end{array}
      \\[1ex]
      &\sim{}
      \left[
        \begin{array}{rrrr|r}
          0 & 0 & 3 & 0 & 0 \\
          0 & 1 & -1 & -1 & 0 \\
          0 & 0 & 1 & 0 & 0 \\
          1 & 0 & 0 & 1 & 0
        \end{array}
      \right]
      \quad
      \begin{array}{@{}l@{}}
        \smash{\rowop{R_1\leftrightarrow R_4}} \\
        {} \\
        {} \\
        {}
      \end{array}
      \\[1ex]
      &\sim{}
      \left[
        \begin{array}{rrrr|r}
          1 & 0 & 0 & 1 & 0 \\
          0 & 1 & -1 & -1 & 0 \\
          0 & 0 & 1 & 0 & 0 \\
          0 & 0 & 3 & 0 & 0
        \end{array}
      \right]
      \quad
      \begin{array}{@{}l@{}}
        {} \\
        \smash{\rowop{+R_3}} \\
        {} \\
        \smash{\rowop{-3R_3}}
      \end{array}
      \\[1ex]
      &\sim{}
      \left[
        \begin{array}{rrrr|r}
          1 & 0 & 0 & 1 & 0 \\
          0 & 1 & 0 & -1 & 0 \\
          0 & 0 & 1 & 0 & 0 \\
          0 & 0 & 0 & 0 & 0
        \end{array}
      \right].
    \end{aligned}
  \]
  Täten \(x_1=-x_4\), \(x_2=x_4\) ja \(x_3=0\), ja muuttuja
  \(x_4\) on vapaa. Valitsemalla \(x_4=t\in\RR\) saadaan
  \[
    x =
    \begin{bmatrix} x_1 \\ x_2 \\ x_3 \\ x_4
    \end{bmatrix}
    =
    \begin{bmatrix} -t \\ t \\ 0 \\ t
    \end{bmatrix}
    = t
    \begin{bmatrix} -1 \\ 1 \\ 0 \\ 1
    \end{bmatrix}.
  \]
  Koska rivioperaatiot säilyttävät ratkaisujoukon, yhtälön \(Ax=0\)
  ratkaisujoukko on
  \[
    \left\{
      t
      \begin{bmatrix} -1 \\ 1 \\ 0 \\ 1
      \end{bmatrix}
      \colon t\in\RR
    \right\}.
  \]
  Esimerkiksi valinnalla \(t=1\) saadaan nollasta poikkeava ratkaisu.
  Lauseen 2.5.1 nojalla \(A\) ei siis ole
  kääntyvä.
\end{solution}

\end{document}
